% Options for packages loaded elsewhere
\PassOptionsToPackage{unicode}{hyperref}
\PassOptionsToPackage{hyphens}{url}
\documentclass[
]{article}
\usepackage{xcolor}
\usepackage{amsmath,amssymb}
\setcounter{secnumdepth}{-\maxdimen} % remove section numbering
\usepackage{iftex}
\ifPDFTeX
  \usepackage[T1]{fontenc}
  \usepackage[utf8]{inputenc}
  \usepackage{textcomp} % provide euro and other symbols
\else % if luatex or xetex
  \usepackage{unicode-math} % this also loads fontspec
  \defaultfontfeatures{Scale=MatchLowercase}
  \defaultfontfeatures[\rmfamily]{Ligatures=TeX,Scale=1}
\fi
\usepackage{lmodern}
\ifPDFTeX\else
  % xetex/luatex font selection
\fi
% Use upquote if available, for straight quotes in verbatim environments
\IfFileExists{upquote.sty}{\usepackage{upquote}}{}
\IfFileExists{microtype.sty}{% use microtype if available
  \usepackage[]{microtype}
  \UseMicrotypeSet[protrusion]{basicmath} % disable protrusion for tt fonts
}{}
\makeatletter
\@ifundefined{KOMAClassName}{% if non-KOMA class
  \IfFileExists{parskip.sty}{%
    \usepackage{parskip}
  }{% else
    \setlength{\parindent}{0pt}
    \setlength{\parskip}{6pt plus 2pt minus 1pt}}
}{% if KOMA class
  \KOMAoptions{parskip=half}}
\makeatother
\usepackage{color}
\usepackage{fancyvrb}
\newcommand{\VerbBar}{|}
\newcommand{\VERB}{\Verb[commandchars=\\\{\}]}
\DefineVerbatimEnvironment{Highlighting}{Verbatim}{commandchars=\\\{\}}
% Add ',fontsize=\small' for more characters per line
\newenvironment{Shaded}{}{}
\newcommand{\AlertTok}[1]{\textcolor[rgb]{1.00,0.00,0.00}{\textbf{#1}}}
\newcommand{\AnnotationTok}[1]{\textcolor[rgb]{0.38,0.63,0.69}{\textbf{\textit{#1}}}}
\newcommand{\AttributeTok}[1]{\textcolor[rgb]{0.49,0.56,0.16}{#1}}
\newcommand{\BaseNTok}[1]{\textcolor[rgb]{0.25,0.63,0.44}{#1}}
\newcommand{\BuiltInTok}[1]{\textcolor[rgb]{0.00,0.50,0.00}{#1}}
\newcommand{\CharTok}[1]{\textcolor[rgb]{0.25,0.44,0.63}{#1}}
\newcommand{\CommentTok}[1]{\textcolor[rgb]{0.38,0.63,0.69}{\textit{#1}}}
\newcommand{\CommentVarTok}[1]{\textcolor[rgb]{0.38,0.63,0.69}{\textbf{\textit{#1}}}}
\newcommand{\ConstantTok}[1]{\textcolor[rgb]{0.53,0.00,0.00}{#1}}
\newcommand{\ControlFlowTok}[1]{\textcolor[rgb]{0.00,0.44,0.13}{\textbf{#1}}}
\newcommand{\DataTypeTok}[1]{\textcolor[rgb]{0.56,0.13,0.00}{#1}}
\newcommand{\DecValTok}[1]{\textcolor[rgb]{0.25,0.63,0.44}{#1}}
\newcommand{\DocumentationTok}[1]{\textcolor[rgb]{0.73,0.13,0.13}{\textit{#1}}}
\newcommand{\ErrorTok}[1]{\textcolor[rgb]{1.00,0.00,0.00}{\textbf{#1}}}
\newcommand{\ExtensionTok}[1]{#1}
\newcommand{\FloatTok}[1]{\textcolor[rgb]{0.25,0.63,0.44}{#1}}
\newcommand{\FunctionTok}[1]{\textcolor[rgb]{0.02,0.16,0.49}{#1}}
\newcommand{\ImportTok}[1]{\textcolor[rgb]{0.00,0.50,0.00}{\textbf{#1}}}
\newcommand{\InformationTok}[1]{\textcolor[rgb]{0.38,0.63,0.69}{\textbf{\textit{#1}}}}
\newcommand{\KeywordTok}[1]{\textcolor[rgb]{0.00,0.44,0.13}{\textbf{#1}}}
\newcommand{\NormalTok}[1]{#1}
\newcommand{\OperatorTok}[1]{\textcolor[rgb]{0.40,0.40,0.40}{#1}}
\newcommand{\OtherTok}[1]{\textcolor[rgb]{0.00,0.44,0.13}{#1}}
\newcommand{\PreprocessorTok}[1]{\textcolor[rgb]{0.74,0.48,0.00}{#1}}
\newcommand{\RegionMarkerTok}[1]{#1}
\newcommand{\SpecialCharTok}[1]{\textcolor[rgb]{0.25,0.44,0.63}{#1}}
\newcommand{\SpecialStringTok}[1]{\textcolor[rgb]{0.73,0.40,0.53}{#1}}
\newcommand{\StringTok}[1]{\textcolor[rgb]{0.25,0.44,0.63}{#1}}
\newcommand{\VariableTok}[1]{\textcolor[rgb]{0.10,0.09,0.49}{#1}}
\newcommand{\VerbatimStringTok}[1]{\textcolor[rgb]{0.25,0.44,0.63}{#1}}
\newcommand{\WarningTok}[1]{\textcolor[rgb]{0.38,0.63,0.69}{\textbf{\textit{#1}}}}
\setlength{\emergencystretch}{3em} % prevent overfull lines
\providecommand{\tightlist}{%
  \setlength{\itemsep}{0pt}\setlength{\parskip}{0pt}}
\usepackage{bookmark}
\IfFileExists{xurl.sty}{\usepackage{xurl}}{} % add URL line breaks if available
\urlstyle{same}
\hypersetup{
  hidelinks,
  pdfcreator={LaTeX via pandoc}}

\author{}
\date{}

\begin{document}

\section{Projection Obstruction for Factorized States: Nonlinear
Dynamics, Additive ℓ¹ Structure, and Finite-Time Defect
Bounds}\label{projection-obstruction-for-factorized-states-nonlinear-dynamics-additive-ux2113uxb9-structure-and-finite-time-defect-bounds}

\textbf{Author:} Jeremy H. Carroll \textbf{Date:} March 2026
\textbf{Version:} v5.0 (Pre-Print)

\begin{center}\rule{0.5\linewidth}{0.5pt}\end{center}

\subsection{Abstract}\label{abstract}

Projection-based approximations --- such as mean-field theory,
product-state ansatze, and factorized variational methods --- replace
correlated states with product states via a nonlinear retraction. This
induces a structural discrepancy between the true dynamics and the
projected dynamics.

We analyze this discrepancy in a Banach presheaf framework and establish
three results. \textbf{(1)} For a \^{}2\$ retraction \(\Pi\) onto a
submanifold \(, the induced dynamics  = \Pi \circ T_t|_M\) is, under a
non-degeneracy condition on the Jacobian of the projected generator, not
conjugate to any linear semigroup (Theorem 3.1). \textbf{(2)} Under
strict locality, additivity, and symmetry axioms (Axioms 1--4), the
\(\ell^1\) norm of the presheaf coboundary is the canonical additive
diagnostic for this discrepancy, characterized as the initial object in
an appropriate category of seminorms (Theorem 4.2). \textbf{(3)} For
finite-range, finite-dimensional lattice systems with non-commuting
interactions, the induced defect exhibits non-vanishing density over a
finite time interval independent of system size, as constrained by
Lieb--Robinson bounds (Theorem 5.1).

Axioms 1--4 characterize a class of local, additive, and
symmetry-respecting defect diagnostics. Within this class, the ℓ¹
coboundary norm is uniquely determined up to a global scalar.

These results formalize the geometric, categorical, and dynamical
structure of projection-induced errors in factorized many-body
approximations. The framework is finite-dimensional and finite-time; no
claims are made regarding asymptotic or thermodynamic limits.

--- such as mean-field theory, product-state ansatze, and factorized
variational methods --- replace correlated states with product states
via a nonlinear retraction. This induces a structural discrepancy
between the true dynamics and the projected dynamics.

We analyze this discrepancy in a Banach presheaf framework and establish
three results. \textbf{(1)} For a \(C^2\) retraction \(\Pi\) onto a
submanifold \(M\), the induced dynamics \(S_t = \Pi \circ T_t|_M\) is,
under a non-degeneracy condition on the Jacobian of the projected
generator, not conjugate to any linear semigroup (Theorem 3.1).
\textbf{(2)} Under strict locality, additivity, and symmetry axioms
(Axioms 1--4), the \(\ell^1\) norm of the presheaf coboundary is the
canonical additive diagnostic for this discrepancy, characterized as the
initial object in an appropriate category of seminorms (Theorem 4.2).
\textbf{(3)} For finite-range, finite-dimensional lattice systems with
non-commuting interactions, the induced defect exhibits non-vanishing
density over a finite time interval independent of system size, as
constrained by Lieb--Robinson bounds (Theorem 5.1).

We further observe that Axioms 1--4 are not arbitrary: they characterize
the minimal structural conditions under which a local agent can
consistently detect, aggregate, and act upon defects (Section 1.2). This
grounds the \(\ell^1\) diagnostic in an operational criterion --- the
existence of consistent local observers --- rather than in unmotivated
mathematical assumptions.

These results formalize the geometric, categorical, and dynamical
structure of projection-induced errors in factorized many-body
approximations. The framework is finite-dimensional and finite-time; no
claims are made regarding asymptotic or thermodynamic limits.

\begin{center}\rule{0.5\linewidth}{0.5pt}\end{center}

\subsection{1. Introduction}\label{introduction}

Projection-based approximations are central to many-body physics and
computation. Methods such as mean-field theory, Hartree--Fock, and
product-state variational ansatze replace a correlated state with a
product of local states, thereby reducing an exponentially large
description to a tractable one. This replacement is implemented by a
nonlinear map --- formally, a retraction --- onto a low-dimensional
manifold of factorized states. While the practical success of these
methods is well established, the structural effect of this retraction on
the underlying dynamics is less clearly understood: what information is
systematically discarded, how it propagates, and how it should be
measured in a representation-independent way.

In this work, we isolate and analyze the discrepancy introduced by
projection as a geometric and combinatorial object. We show that
composing a linear many-body evolution with a nonlinear retraction
produces dynamics that --- under a non-degeneracy condition --- cannot
be reduced to any linear flow, and that the resulting local
inconsistencies assemble into a presheaf coboundary with a canonical
additive structure. Under natural locality and symmetry conditions, this
leads to a distinguished \(\ell^1\) diagnostic for projection-induced
defects, and for finite-range systems we prove that this defect persists
with non-vanishing density over a finite time interval independent of
system size.

The goal is not to propose a new physical theory, but to provide a
precise structural account of the error introduced by factorization.

We prove three structural theorems: 1. \textbf{Geometric Obstruction
(Theorem 3.1):} The projected dynamics is not \(C^2\)-conjugate to any
linear semigroup when the projected generator has non-constant Jacobian.
2. \textbf{Additive Universality (Theorem 4.2):} Subject to Axioms 1--4,
the \(\ell^1\) coboundary norm is the uniquely determined diagnostic, up
to a global scalar. Relaxing these axioms admits alternative norms
(e.g.~\(\ell^2\)). 3. \textbf{Density Persistence (Theorem 5.1):} The
per-bond obstruction density is \(\Omega(1)\) in system size over a time
window \(O(1/J)\), with constants depending only on local dimension,
interaction strength, and active bond fraction.

\textbf{Scope.} The framework applies rigorously to Banach presheaves
over finite posets. We avoid speculative physics or infinite-time
limits; the density scaling bound is finite-time only.

\textbf{Conventions.} Norms are denoted \(\lVert \cdot \rVert\) with
appropriate subscripts; scalar absolute values are denoted \(|\cdot|\).

\subsubsection{1.1 Relation to Companion
Work}\label{relation-to-companion-work}

The additive diagnostic used throughout this paper --- the \(\ell^1\)
coboundary norm --- is not introduced ad hoc. It is the terminal output
of a two-stage classification program establishing \(\ell^1\) rigidity
at progressively richer structural levels:

\begin{itemize}
\tightlist
\item
  \textbf{Stage 1 (combinatorial):} In ``Coordinate-Wise Additivity and
  the \(\ell^1\) Norm on Finite Graph Cochains'' {[}8{]}, local
  additivity over disjoint coordinate supports and edge separability on
  finite graph cochains already force an \(\ell^1\) geometry. No Banach
  space structure or dynamics are required.
\item
  \textbf{Stage 2 (functional-analytic):} In ``The Free \(\ell^1\)
  Seminorm on Banach Presheaf Coboundaries'' {[}6{]}, the same rigidity
  persists when the cochain spaces are replaced by Banach presheaves and
  coordinate additivity is replaced by structural compatibility
  (functoriality under bounded linear maps).
\end{itemize}

By the classification results of these preceding papers, any diagnostic
satisfying locality, additivity over disjoint supports, and structural
compatibility with the underlying linear maps must be proportional to
the \(\ell^1\) coboundary norm. Accordingly, the defect quantity
analyzed here is not chosen, but structurally determined by these
constraints.

The present paper (Stage 3, dynamical) builds on that classification by
showing that projection-induced dynamics naturally generate nontrivial
coboundary defects, and that these defects persist with non-vanishing
density under finite-range evolution. Together, the three papers
establish that additive, structure-preserving measurements of local
defects are uniquely forced to be \(\ell^1\) across combinatorial,
functional, and dynamical settings, and that projection dynamics
generically produce nonzero instances of such defects.

\subsubsection{Conceptual Overview}\label{conceptual-overview}

\begin{Shaded}
\begin{Highlighting}[]
\NormalTok{Full correlated state space X}
\NormalTok{        |}
\NormalTok{        |  C\^{}2 retraction Pi}
\NormalTok{        v}
\NormalTok{Product{-}state manifold M}
\NormalTok{        |}
\NormalTok{        |  induced curvature term D\^{}2Pi(Am, .)}
\NormalTok{        v}
\NormalTok{Projected dynamics S\_t = Pi . T\_t|\_M}
\NormalTok{        |}
\NormalTok{        |  local incompatibilities across regions}
\NormalTok{        v}
\NormalTok{Coboundary defects delta\_e(s)}
\NormalTok{        |}
\NormalTok{        |  additive aggregation (Axioms 1{-}4)}
\NormalTok{        v}
\NormalTok{Total obstruction  I(s) = Sum ||delta\_e(s)||}
\end{Highlighting}
\end{Shaded}

\begin{center}\rule{0.5\linewidth}{0.5pt}\end{center}

\subsection{2. Structural Framework}\label{structural-framework}

We work with Banach-valued presheaves over finite posets, following
standard categorical conventions {[}2{]}.

\subsubsection{2.1 Presheaves over Finite
Posets}\label{presheaves-over-finite-posets}

Let \((P, \leq)\) be a finite poset with covering relations
\(E_{\text{cov}}(P) = \{(a, b) : a < b,\; \nexists\, c \text{ with } a < c < b\}\).
A \textbf{Banach presheaf} \(F\) over \(P\) assigns to each \(x \in P\)
a Banach space \(F(x)\) and to each \((a,b) \in E_{\text{cov}}\) a
contractive linear restriction map \(r_{b,a}: F(b) \to F(a)\).

\subsubsection{2.2 The Coboundary}\label{the-coboundary}

For each edge \(e = (a,b)\), the \textbf{coboundary} is
\(\delta_e(s) = r_{b,a}(s(b)) - s(a) \in F(a)\). A section descends if
\(\delta_e(s) = 0\) for all \(e\). The \(\ell^1\) coboundary norm is
\(I(s) = \sum_{e \in E_{\text{cov}}} \|\delta_e(s)\|\).

\subsubsection{2.3 Linear Semigroups and Projected
Dynamics}\label{linear-semigroups-and-projected-dynamics}

A \textbf{\(C_0\) contraction semigroup} on a Banach space \(X\) is
\(\{T_t\}_{t \geq 0}\) with generator \((A, \mathcal{D}(A))\). A
\textbf{\(C^2\) retraction} is a \(C^2\) idempotent
\(\Pi: X \to M \subset X\) with \(M\) a regular \(C^2\) submanifold. If
\(\{T_t\}\) does not preserve \(M\), we define the \textbf{projected
semigroup} \(S_t = \Pi \circ T_t|_M\) mapping \(M \to M\) locally.

\begin{center}\rule{0.5\linewidth}{0.5pt}\end{center}

\subsection{3. Theorem I: Nonlinear Obstruction to Linear
Conjugacy}\label{theorem-i-nonlinear-obstruction-to-linear-conjugacy}

\subsubsection{3.1 Statement}\label{statement}

\begin{quote}
\textbf{Theorem 3.1 (Projection Non-Dilability).} Let \(X\) be a Banach
space, \(\{T_t\}\) a \(C_0\) contraction semigroup with generator
\((A, \mathcal{D}(A))\), and \(\Pi: X \to M \subset X\) a \(C^2\)
retraction onto a \(C^2\) submanifold \(M \subset \mathcal{D}(A)\).
Suppose:

\begin{enumerate}
\def\labelenumi{(\alph{enumi})}
\item
  \emph{(Non-invariance)} There exists \(t > 0\) such that
  \(T_t(M) \not\subset M\).
\item
  \emph{(Non-degeneracy)} The induced vector field \(B(m) = D\Pi(m)Am\)
  has non-constant Jacobian \(DB(m)\) in any affine coordinate chart
  compatible with a flat connection.
\end{enumerate}

Then there does not exist a Banach space \(Y\), a \(C^2\) embedding
\(\iota: M \hookrightarrow Y\), and a \(C_0\) contraction semigroup
\(\{\tilde{T}_t\}\) on \(Y\) such that
\(S_t = \iota^{-1} \circ \tilde{T}_t \circ \iota\). All arguments are
local in affine coordinate charts induced by the ambient Banach space;
global flatness of \(M\) is not required.
\end{quote}

\subsubsection{3.2 Preliminary Lemma}\label{preliminary-lemma}

\begin{quote}
\textbf{Lemma 3.2 (Parallel fields on flat connections).} Let
\((M, \nabla)\) be a connected Banach manifold equipped with a flat
affine connection (\(R^{\nabla} = 0\)). If a \((1,1)\)-tensor field
\(S\) on \(M\) satisfies \(\nabla S = 0\), then \(S\) is constant in any
affine coordinate chart.
\end{quote}

\emph{Proof.} Since \(R^{\nabla} = 0\), the exponential map provides
local affine coordinates where the Christoffel symbols vanish. In these
coordinates, \(\nabla S = 0\) reduces to \(\partial_i S^j_k = 0\),
making components constant. Affine transitions preserve this constancy
globally. \(\square\)

\subsubsection{3.3 Proof of Theorem 3.1}\label{proof-of-theorem-3.1}

\textbf{Step 1: Linear flows are affinely flat.} A linear semigroup
\(\{\tilde{T}_t\}\) with generator \(\tilde{A}\) satisfies
\(\tilde{A}(y) = \tilde{A} y\). The canonical flat connection
\(\tilde{\nabla}\) on \(Y\) has \(R^{\tilde{\nabla}} = 0\), and
\(\tilde{\nabla}^2 \tilde{A} = 0\).

\textbf{Step 2: Conjugacy transfers flatness.} Suppose
\(\iota: M \hookrightarrow Y\) is a \(C^2\) embedding conjugating
\(S_t\) to \(\tilde{T}_t\). The pullback connection
\(\nabla = \iota^* \tilde{\nabla}\) is flat on \(M\) (since \(\iota\) is
\(C^2\), the pullback preserves the vanishing of curvature), and
\(\nabla^2 B = 0\) for the generator \(B = \iota^{-1}_* \tilde{A}\) of
\(S_t\).

\textbf{Step 3: The projected generator has non-constant derivative.}
The projected generator is \(B(m) = D\Pi(m)Am\). Differentiating yields:
\[\nabla_v B = D^2\Pi(m)(Am, v) + D\Pi(m)Av\]

If \(\nabla^2 B = 0\), then \(\nabla_v B\) would be a parallel tensor
field on \((M, \nabla)\), and by Lemma 3.2, it would be constant in any
affine coordinate chart. But Hypothesis (b) states that \(DB(m)\) is
non-constant in every such chart. This is a contradiction. \(\square\)

\begin{center}\rule{0.5\linewidth}{0.5pt}\end{center}

\subsection{\texorpdfstring{4. Theorem II: Additive Universality of the
\(\ell^1\)
Diagnostic}{4. Theorem II: Additive Universality of the \textbackslash ell\^{}1 Diagnostic}}\label{theorem-ii-additive-universality-of-the-ell1-diagnostic}

Theorem 3.1 demonstrates that geometric nonlinearity is inevitably
produced by the retraction. The next question is how to measure the
resulting defect. Theorem 4.2 establishes that, given Axioms 1--4, the
\(\ell^1\) norm is the uniquely forced diagnostic. The conclusion is
conditional on these axioms; relaxing them admits other norms.

\subsubsection{4.1 Diagnostic Systems and the Category of
Seminorms}\label{diagnostic-systems-and-the-category-of-seminorms}

Let \(C^1(P, F) = \bigoplus_{e \in E_{\text{cov}}} F(a_e)\) be the space
of edge defects. We consider seminorms
\(J: C^1(P,F) \to \mathbb{R}_{\geq 0}\) satisfying:

\textbf{Axiom 1 (Coboundary Locality).}
\(J_P(s) = \sum_{e} w_e \|\delta_e s\|\), with local weights
\(w_e > 0\).

\textbf{Axiom 2 (Faithfulness).} \(J_P(s) = 0 \iff s \text{ descends}\).

\textbf{Axiom 3 (Contractive Monotonicity).} For local contractions, the
norm weakly decreases.

\textbf{Axiom 4 (Symmetry and Coproduct Compatibility).} Edge weights
are uniform by edge-exchange symmetry, and the norm splits linearly over
disjoint graph unions.

\subsubsection{4.2 Statement}\label{statement-1}

\begin{quote}
\textbf{Theorem 4.2 (\(\ell^1\) Additive Universality).} The \(\ell^1\)
coboundary norm \(I(s) = \sum_e \|\delta_e(s)\|\) is the initial object
in the category of faithful seminorms on \(C^1(P,F)\) satisfying Axioms
1--4 (with morphisms given by monotone seminorm contractions), in the
sense that any other such seminorm factors uniquely through \(I\) up to
a positive scalar. Consequently, any diagnostic satisfying these axioms
is proportional to \(I\): \[J(s) = c \sum_e \|\delta_e(s)\|\] for some
constant \(c > 0\).
\end{quote}

\emph{Proof Sketch.} Any diagnostic \(J_P\) satisfying the axioms must
assign constant weights to local defects due to edge-exchange
isomorphism (Axiom 4). Concretely, edge-exchange symmetry implies that
all edge weights \(w_e\) are equal, since any permutation of edges that
preserves graph structure must leave \(J\) invariant. The Kan extension
\(\mathrm{Lan}_\pi d(s)\) realizes the universal space of local defects.
Additivity over disjoint unions (Axiom 4) forces the seminorm to
decompose as a sum over edges, excluding \(\ell^2\) or \(\ell^\infty\)
alternatives, which do not decompose additively over disjoint unions.
Thus \(\ell^1\) is uniquely selected. \(\square\)

\emph{Remark.} The \(\ell^1\) norm is uniquely selected \textbf{given
Axioms 1--4}, which may or may not be physically appropriate depending
on context. The structural content of the theorem is that \(\ell^1\)
corresponds to additive aggregation of independent defects. In settings
where defect channels are not independent, alternative diagnostics may
be appropriate.

By the classification results of the preceding papers {[}8, 6{]} ---
first at the combinatorial level (local additivity over disjoint
supports) and then at the Banach level (functoriality under bounded
linear maps) --- any diagnostic satisfying locality, additivity, and
structural compatibility must be proportional to the \(\ell^1\)
coboundary norm. Accordingly, the defect quantity analyzed here is not
chosen, but structurally determined by these constraints.

\begin{center}\rule{0.5\linewidth}{0.5pt}\end{center}

\subsection{5. Theorem III: Finite-Time Density
Persistence}\label{theorem-iii-finite-time-density-persistence}

\subsubsection{5.1 Setup}\label{setup}

Consider \(N\) sites on a finite graph \(\Lambda\) with local Hilbert
spaces (\(\dim \mathcal{H}_x \leq d < \infty\)), bounded Hamiltonian
\(H = \sum_e H_e\) with \(\|H_e\|_{\text{op}} \leq J\), and retraction
\(\Pi(\rho) = \bigotimes_x \text{Tr}_{\bar{x}}(\rho)\). Local
Hamiltonians satisfying the non-commutativity requirement below are
generic in the space of two-local interactions {[}3{]}.

Define the commutator:
\[[H_e, \rho_a \otimes \rho_b] = H_e(\rho_a \otimes \rho_b) - (\rho_a \otimes \rho_b)H_e\]

\subsubsection{5.2 Explicit Rate Bound}\label{explicit-rate-bound}

\begin{quote}
\textbf{Lemma 5.2 (Local obstruction rate).} For a single bond
\(e = (a,b)\) with \([H_e, \rho_a \otimes \rho_b] \neq 0\), let
\(\gamma_e = \|[H_e, \rho_a \otimes \rho_b]\|_1 / d^2\). Then
\(\gamma_e > 0\). For \(t\) small:
\[C_e(t) \geq \gamma_e t - O(J^2 t^2).\]
\end{quote}

\emph{Proof.} On finite-dimensional spaces, the map
\(t \mapsto \rho(t) - \Pi(\rho(t))\) is analytic. Expanding to first
order, the leading term is \(-it[H_e, \rho_a \otimes \rho_b]\), whose
trace norm is \(\gamma_e d^2 / d^2 = \gamma_e\). The remainder is
bounded by \(O(J^2 t^2)\) via the operator norm bound
\(\|H_e\|_{\text{op}} \leq J\). Since \(\gamma_e > 0\) by hypothesis,
the linear term dominates for \(t < \gamma_e / (KJ^2)\) where \(K\) is a
computable constant depending only on \(d\). \(\square\)

\subsubsection{5.3 The Density Bound}\label{the-density-bound}

\begin{quote}
\textbf{Theorem 5.1 (Finite-Time Density Persistence).} Let \(\Lambda\)
be a finite graph with \(N\) sites and finite-dimensional local Hilbert
spaces (\(\dim \leq d\)). Let \(H = \sum_e H_e\) be a bounded
finite-range Hamiltonian with \(\|H_e\|_{\text{op}} \leq J\), and
suppose a fraction \(\eta > 0\) of edges satisfy
\([H_e, \rho_a \otimes \rho_b] \neq 0\) for an initial product state
\(\rho(0) = \bigotimes_x \rho_x\).

Then there exist constants \(c > 0\) and \(t_0 > 0\), depending only on
local dimension \(d\), interaction strength \(J\), and active fraction
\(\eta\) (but not on system size \(N\)), such that
\[\frac{C(t)}{N} \geq c \quad \text{for all } t \in [t_1, t_2] \subset (0, t_0]\]
where \(t_0 = O(1/J)\).
\end{quote}

\emph{Proof sketch.} By Lemma 5.2, each active bond \(e\) generates a
trace-norm defect \(C_e(t) \geq \gamma_e t - O(J^2 t^2)\) for small
\(t\). By the Lieb--Robinson bound {[}1, 4, 5{]}, for times
\(t < R / v_{\text{LR}}\) (where \(R\) is the interaction range and
\(v_{\text{LR}}\) the Lieb--Robinson velocity), the reduced density
matrix on any bond is approximated by its isolated evolution up to an
error exponentially small in \((v_{\text{LR}} t - R)\). Concretely, the
trace-norm error from truncating the causal cone satisfies
\(\|\rho_e^{\text{exact}}(t) - \rho_e^{\text{local}}(t)\|_1 \leq C_{\text{LR}} e^{-\mu(R - v_{\text{LR}} t)}\)
for constants \(C_{\text{LR}}, \mu > 0\) depending on the lattice
geometry.

Within this causal window, each bond's defect \(C_e(t)\) is controlled
by the local evolution and cannot be canceled by correlations from
distant bonds, because such correlations have not yet propagated to the
bond. Summing over the \(\eta N\) active bonds and noting that
\(\gamma_e \geq \gamma_{\min} > 0\) (where \(\gamma_{\min}\) depends on
\(J\) and \(d\)), we obtain
\(C(t) \geq \eta N (\gamma_{\min} t - KJ^2 t^2)\) for
\(t \in (0, \gamma_{\min}/(2KJ^2))\). Dividing by \(N\) gives
\(C(t)/N \geq \eta \gamma_{\min} t / 2 > 0\) on a non-trivial interval,
with all constants independent of \(N\). \(\square\)

\emph{Remark on the time window.} The interval \([t_1, t_2]\) scales as
\(O(1/J)\): the defect appears immediately (linear growth at rate
\(\gamma_e\)), dominates the quadratic correction for \(t \ll 1/J\), and
the Lieb--Robinson causal cone remains localized for
\(t \ll R/v_{\text{LR}} \sim 1/J\). Beyond this window, many-body
correlations propagate across bonds and the argument does not apply.

\begin{center}\rule{0.5\linewidth}{0.5pt}\end{center}

\subsection{6. Worked Examples}\label{worked-examples}

\subsubsection{6.1 Example 1: Two Qubits
(Entangling)}\label{example-1-two-qubits-entangling}

Consider two qubits with \(H = J(X \otimes X)\) and initial state
\(|{00}\rangle\). The commutator \([H, |00\rangle\langle00|]\) is
strictly non-zero. The full evolution yields:
\[|\psi(t)\rangle = \cos(Jt)|{00}\rangle - i\sin(Jt)|{11}\rangle\]
\[\rho(t) = \cos^2(Jt)|{00}\rangle\langle{00}| + \sin^2(Jt)|{11}\rangle\langle{11}| + (\text{off-diagonal terms})\]

The product retraction isolates marginals:
\[\rho_A(t) = \rho_B(t) = \cos^2(Jt)|{0}\rangle\langle{0}| + \sin^2(Jt)|{1}\rangle\langle{1}|\]
\[(\Pi\rho)(t) = \cos^4(Jt)|{00}\rangle\langle{00}| + \sin^4(Jt)|{11}\rangle\langle{11}| + \cos^2(Jt)\sin^2(Jt)(|{01}\rangle\langle{01}| + |{10}\rangle\langle{10}|)\]

The trace norm defect \(C(t) = \|\rho - \Pi\rho\|_1 > 0\) for \(t > 0\).
This witnesses non-product structure in \(\rho(t)\): the true state
cannot be written as a tensor product of marginals, and the defect
quantifies this failure.

\subsubsection{6.2 Example 2: Transverse-Field Ising
Model}\label{example-2-transverse-field-ising-model}

The 1D Ising model \(H = -J\sum_i Z_i Z_{i+1} - h\sum_i X_i\) with
\(|{+}\rangle^{\otimes N}\) produces strictly positive initial rates
\(C_e(t) = \gamma_e t + O(J^2 t^2) > 0\), consistent with the
linear-in-\(t\) scaling established by Lemma 5.2. By translational
invariance, per-bond density is insensitive to \(N\) for
\(t \ll N/v_{\text{LR}}\).

\begin{center}\rule{0.5\linewidth}{0.5pt}\end{center}

\subsection{7. Discussion}\label{discussion}

This work isolates a structural feature common to projection-based
approximations: the act of restricting a linear many-body evolution to a
nonlinear manifold introduces an intrinsic obstruction that is both
geometrically unavoidable and locally detectable. The three results
established here --- nonlinear non-dilability of the projected flow,
additive universality of the \(\ell^1\) coboundary diagnostic under
symmetry and locality constraints, and finite-time persistence of defect
density in finite-range systems --- together provide a coherent
description of how information is lost under factorization and how that
loss propagates. Importantly, these statements are finite-dimensional
and finite-time, and rely only on structural properties of the
retraction and locality of interactions, rather than on model-specific
features. In this sense, the framework does not compete with existing
approximation methods, but rather clarifies the limits they necessarily
obey.

Looking ahead, just as differential forms decompose into local and
global topological pieces, resolving \(C^1(P, F)\) defects suggests a
generalized \textbf{Hodge decomposition} for dynamical lattice models:
\[ \delta s = d\alpha + \delta\beta + h \] where \(d\alpha\) isolates
removable trivial projection errors, \(\delta\beta\) captures dynamical
correlation flux, and \(h\) localizes persistent obstruction invariants
(harmonic forms). This decomposition, together with its duality
structure, is developed in companion papers {[}6, 7{]}.

\begin{center}\rule{0.5\linewidth}{0.5pt}\end{center}

\subsection{Acknowledgments}\label{acknowledgments}

No external funding. No conflicts of interest.

\begin{center}\rule{0.5\linewidth}{0.5pt}\end{center}

\subsection{References}\label{references}

\begin{enumerate}
\def\labelenumi{\arabic{enumi}.}
\tightlist
\item
  E. H. Lieb and D. W. Robinson, ``The Finite Group Velocity of Quantum
  Spin Systems,'' \emph{Commun. Math. Phys.} \textbf{28} (1972),
  251--257.
\item
  S. Mac Lane, \emph{Categories for the Working Mathematician},
  Springer, 1971.
\item
  J. Klassen and B. M. Terhal, ``Two-local qubit Hamiltonians: when are
  they stoquastic?,'' \emph{Quantum} \textbf{3} (2019), 139.
\item
  B. Nachtergaele and R. Sims, ``Lieb--Robinson Bounds and the
  Exponential Clustering Theorem,'' \emph{Commun. Math. Phys.}
  \textbf{265} (2006), 119--130.
\item
  S. Bravyi, M. B. Hastings, and F. Verstraete, ``Lieb--Robinson Bounds
  and the Generation of Correlations and Topological Quantum Order,''
  \emph{Phys. Rev.~Lett.} \textbf{97} (2006), 050401.
\item
  J. H. Carroll, ``The Free \(\ell^1\) Seminorm on Banach Presheaf
  Coboundaries,'' Zenodo, 2026. https://doi.org/10.5281/zenodo.18897080
\item
  J. H. Carroll, ``Hodge Structure and Gauge Equivalence of \(\ell^1\)
  Defect Fields,'' Zenodo, 2026.
\item
  J. H. Carroll, ``Coordinate-Wise Additivity and the \(\ell^1\) Norm on
  Finite Graph Cochains,'' Zenodo, 2026.
  https://doi.org/10.5281/zenodo.18945342
\end{enumerate}

\end{document}
